\section{引言}
\label{sec:introduction}


% 1
coding agent 正被越来越多地部署到长程软件工程任务上，在真实代码仓库的问题（issue）解决~\citep{jimenezSWEbenchCanLanguage2023,yangSWEbenchMultimodalAI2024,dengSWEBenchProCan2025}以及多步终端工作流~\citep{merrillTerminalBenchBenchmarkingAgents2026a}上取得了可度量的进展。
在实践中，这类进展不仅依赖底层的语言模型，同样依赖围绕其构建的工程组件：塑造工作风格的系统提示词、把文件系统与 shell 暴露出来的工具，以及控制上下文、执行与恢复的中间件。
这一组位于模型之外、可编辑的组件被统称为智能体的 \emph{harness}~\citep{rajasekaranHarnessDesignLongrunning2026,lopopoloHarnessEngineeringLeveraging2026,wangOpenHandsOpenPlatform2025,yangSWEagentAgentComputerInterfaces2024,steinbergerOpenClawPersonalAI2026,HermesAgentAgent}。

% 2

即使固定基座模型，harness 设计也会实质性地改变长程编码基准上的任务完成情况~\citep{trivedyImprovingDeepAgents2026,wangOpenHandsOpenPlatform2025}，这使 harness 工程成为改进 coding agent 的一级杠杆。
而且，最优的 harness 是模型相关的：为某一基座模型调优的 harness 在另一模型上往往表现不佳，必须随着基座模型的更换而重新适配。
在当前实践中，这种适配由人工完成——开发者检查轨迹、识别反复出现的失败模式，并跨提示词、工具、中间件与技能手工编写修改。
然而，随着基座模型快速演进~\citep{xiaomimimoteamMiMoV25Pro2026,teamQwen36PlusRealWorld2026,yangQwen3TechnicalReport2025a,DeepSeek_V4pdfDeepseekaiDeepSeekV4Pro2026,kimiteamKimiK26Tech2026,teamKimiK25Visual2026}，这一人工闭环难以跟上节奏，在模型能力与释放该能力所需的 harness 之间造成了不断扩大的差距~\citep{steinbergerOpenClawPersonalAI2026}。

% 3

一个直观的方向是用一个\textit{演化智能体}来自动化这一循环，由它基于经验来优化 harness 组件~\citep{agrawalGEPAReflectivePrompt2025a,zhangAgenticContextEngineering2025a,caiTrainingFreeGroupRelative2025}。
然而，现有方法中鲜有能联合演化全部可编辑组件的~\citep{leeMetaHarnessEndtoEndOptimization2026}；大多数只聚焦于单一组件，通常是提示词~\citep{shinnReflexionLanguageAgents2023c,zhaoExpeLLLMAgents2024,madaanSelfRefineIterativeRefinement2023c}、技能~\citep{maSkillClawLetSkills2026b,xiaSkillRLEvolvingAgents2026c}或上下文中的操作手册（playbook）~\citep{zhangAgenticContextEngineering2025a}。
端到端地联合演化多个组件面临两个结构性障碍：冗长且无结构的轨迹几乎提供不了可付诸行动的信号；而高度耦合的 harness 框架使得提示词之外的修改极易出错。
这使得智能体驱动的 harness 演化的核心问题仍然悬而未决：
\textit{演化智能体如何才能\textbf{联合地}且\textbf{稳定地}演化 coding agent harness 的全部可编辑组件？}

% 4

我们的核心洞见是：该问题的瓶颈在于\emph{可观测性}，而不在于智能体的能力——一旦演化智能体在一个清晰的动作空间之上获得了结构化的上下文，它就能可靠地收敛到更好的 harness 设计~\citep{suttonBitterLesson2019,zunicBitterLessonAgent2026}。
我们在 \textbf{A}gentic \textbf{H}arness \textbf{E}ngineering
（\textbf{AHE}，图~\ref{fig:method}）中实现了这一点，它是一个由三大可观测性支柱驱动的闭环：
\ding{182}~\emph{组件可观测性}，通过一个解耦的 harness 把七类可编辑组件以文件形式暴露出来，使每一种失败模式都能干净地映射到单一的组件类别；
\ding{183}~\emph{经验可观测性}，通过从数百万原始轨迹 token 中蒸馏出的分层、可逐级下钻的证据语料，使演化者消费的是结构化的根因而非原始日志；以及
\ding{184}~\emph{决策可观测性}，通过一份变更清单为每一次修改配上自我声明的预测，随后用下一轮的任务级结果对其加以验证，从而使每一次修改都成为一份可证伪的契约，无效的修改则以文件粒度被回滚。

% 5

我们在 Terminal-Bench 2\citep{merrillTerminalBenchBenchmarkingAgents2026a}上对 AHE 进行了实证验证：十轮迭代把 pass@1 从 69.7\% 提升到 77.0\%，超过了人工设计的 Codex~\citep{openaiCodexCLI2025}以及自演化基线 ACE~\citep{zhangAgenticContextEngineering2025a}与 Training-Free GRPO~\citep{caiTrainingFreeGroupRelative2025}。
% Without further evolution, the frozen harness transfers to SWE-bench-verified~\citep{jimenezSWEbenchCanLanguage2023} and to three alternate base-model families, with the largest gains on bases further from saturation, suggesting that AHE encodes coordination patterns that less-saturated models lean on more heavily.
无需进一步演化，冻结后的 harness 即可迁移到 SWE-bench-verified~\citep{jimenezSWEbenchCanLanguage2023}；在另外三个基座模型家族上，它带来了 $+5.1$ 到 $+10.1$ 个百分点的一致 pass@1 提升，其中距离饱和更远的基座模型收益最大，这表明 AHE 编码了越是未饱和的模型越依赖的协调模式。
组件消融进一步定位了这一增益的来源：工具、中间件与长期记忆各自单独都能承载这一提升，而仅有系统提示词时反而出现回归，这说明事实性的 harness 结构能够跨任务、跨模型迁移，而散文层面的策略则不能。

本文做出三点贡献：
\begin{itemize}
    \item 我们为 coding agent 形式化了\emph{智能体驱动的 harness 演化}问题，并提出 \textbf{AHE}：它将\emph{跨组件、跨轨迹与跨决策的可观测性}确立为设计枢纽，并通过三大可观测性支柱把每一次 harness 修改都变成可证伪的、文件级的契约——一个解耦的组件基底、一条分层的轨迹蒸馏流水线，以及一份其自我声明的预测由下一轮任务差值加以验证的变更清单。
    \item 我们通过实验表明，AHE 把 Terminal-Bench~2 上的 pass@1 从 69.7\% 提升到 77.0\%，超越了人工设计的与自动化的基线，并产出了一个能够跨基准、跨基座模型家族迁移的冻结 harness。
    \item 我们的分析揭示了智能体驱动演化的两个局限：harness 组件之间以非可加的方式相互作用，因此叠加多个有效修改会使总体增益封顶；以及该闭环的自归因对于修复是可靠的，却对回归视而不见，这指明了回归的前瞻预判是未来自演化闭环最清晰的改进方向。
\end{itemize}

